\documentclass[10pt]{scrartcl}
\usepackage[mdthm,colorsec,noauthor,noasy]{evan}
\usepackage{handout}
\newcommand{\aug}[2]{\left[\begin{array}{@{}#1@{}}#2\end{array}\right]}
\newcommand{\vv}[1]{\mathbf{#1}}
\DeclareMathOperator{\spn}{Span}\DeclareMathOperator{\Nul}{Nul}\DeclareMathOperator{\Col}{Col}\DeclareMathOperator{\Row}{Row}\DeclareMathOperator{\tr}{tr}
\title{Math 54 --- Week 7}
\subtitle{Solutions to the problem set}
\author{}\date{}\ihead{}
\begin{document}
\maketitle

\begin{prob}\textbf{\xlink{math54}{2026-10-05-coordinates-eigenvalues}{prob-P1}{P1}.}\par\idea{Solve $c_1\vv b_1+c_2\vv b_2=\vv x$ by row reducing $[\,\vv b_1\ \vv b_2\mid\vv x\,]$.}
\begin{steps}
  \ii Row reduce:
  \[
    \aug{cc|c}{1&2&4\\3&5&9} \xrightarrow{R_2-3R_1} \aug{cc|c}{1&2&4\\0&-1&-3}.
  \]
  \ii Back-substitute: $c_2=3$, then $c_1=4-2c_2=-2$. So $[\vv x]_{\mathcal B}=(-2,3)$.
  \ii Check: $-2\begin{psmallmatrix}1\\3\end{psmallmatrix}+3\begin{psmallmatrix}2\\5\end{psmallmatrix}=\begin{psmallmatrix}-2+6\\-6+15\end{psmallmatrix}=\begin{psmallmatrix}4\\9\end{psmallmatrix}$.
\end{steps}
\end{prob}

\begin{prob}\textbf{\xlink{math54}{2026-10-05-coordinates-eigenvalues}{prob-P2}{P2}.}\par\idea{Work with coefficient vectors in the standard basis $\{1,t,t^2\}$; coordinates keep linear relations (\xlink{math54}{2026-10-05-coordinates-eigenvalues}{thm-lay-4-9}{Thm.~4.4.9}).}
\begin{steps}
  \ii \textbf{(a)} Solve $c_1(1+t)+c_2(t+t^2)+c_3(1+t^2)=3+2t-t^2$. Matching coefficients of $1,t,t^2$:
  \[
    c_1+c_3=3,\qquad c_1+c_2=2,\qquad c_2+c_3=-1.
  \]
  Adding all three gives $2(c_1+c_2+c_3)=4 \implies c_1+c_2+c_3=2$. Subtracting each gives $c_2=2-3=-1$, $c_3=2-2=0$, $c_1=2-(-1)=3$. Thus $[p]_{\mathcal B}=(3,-1,0)$. Check: $3(1+t)-(t+t^2)=3+2t-t^2$.
  \ii \textbf{(b)} Coordinate vectors in $\{1,t,t^2\}$ are $(1,-1,0)$, $(0,1,-1)$, $(1,0,-1)$. Since $(1,0,-1)=(1,-1,0)+(0,1,-1)$, we have $1-t^2=(1-t)+(t-t^2)$, so the polynomials are \textbf{linearly dependent}.
\end{steps}
\end{prob}

\begin{prob}\textbf{\xlink{math54}{2026-10-05-coordinates-eigenvalues}{prob-P3}{P3}.}\par\idea{Row reduce $[\,\mathcal C\mid\mathcal B\,]$ to $[\,I\mid P_{\mathcal C\leftarrow\mathcal B}\,]$; invert for $P_{\mathcal B\leftarrow\mathcal C}$.}
\begin{steps}
  \ii Row reduce:
  \[
    \aug{cc|cc}{1&0&1&3\\1&1&2&4}
    \xrightarrow{R_2-R_1}
    \aug{cc|cc}{1&0&1&3\\0&1&1&1},
    \qquad
    P_{\mathcal C\leftarrow\mathcal B}=\begin{bmatrix}1&3\\1&1\end{bmatrix}.
  \]
  \ii Determinant is $1-3=-2$, so by the $2\times 2$ inverse formula:
  \[
    P_{\mathcal B\leftarrow\mathcal C}=\frac{1}{-2}\begin{bmatrix}1&-3\\-1&1\end{bmatrix}=\begin{bmatrix}-1/2&3/2\\1/2&-1/2\end{bmatrix}.
  \]
  \ii Then $[\vv x]_{\mathcal C}=P_{\mathcal C\leftarrow\mathcal B}[\vv x]_{\mathcal B}=\begin{psmallmatrix}1&3\\1&1\end{psmallmatrix}\begin{psmallmatrix}2\\-1\end{psmallmatrix}=\begin{psmallmatrix}-1\\1\end{psmallmatrix}$.
  \ii Check: $\vv x=2(1,2)-(3,4)=(-1,0)$, and $-1(1,1)+1(0,1)=(-1,0)$.
\end{steps}
\end{prob}

\begin{prob}\textbf{\xlink{math54}{2026-10-05-coordinates-eigenvalues}{prob-P4}{P4}.}\par\idea{Compute $A\vv v$; for $\lambda=4$, test invertibility of $A-4I$.}
\begin{steps}
  \ii Multiplying directly:
  \[
    A\begin{bmatrix}1\\-1\\0\end{bmatrix} = \begin{bmatrix}3+0+0\\-1-2+0\\0\end{bmatrix} = \begin{bmatrix}3\\-3\\0\end{bmatrix} = 3\begin{bmatrix}1\\-1\\0\end{bmatrix},
  \]
  so $\vv v$ \textbf{is} an eigenvector for eigenvalue $\lambda=3$.
  \ii Expanding $\det(A-4I)$ down column~$2$:
  \[
    A-4I=\begin{bmatrix}-1&0&1\\-1&-2&-1\\0&0&-3\end{bmatrix},\qquad
    \det(A-4I)=-2\det\begin{bmatrix}-1&1\\0&-3\end{bmatrix}=-2(3)=-6\ne 0.
  \]
  Hence $A-4I$ is invertible, so $4$ is \textbf{not} an eigenvalue.
\end{steps}
\end{prob}

\begin{prob}\textbf{\xlink{math54}{2026-10-05-coordinates-eigenvalues}{prob-P5}{P5}.}\par\idea{Column~$2$ of $A-\lambda I$ has only one nonzero entry; expand down it and keep factors.}
\begin{steps}
  \ii Expanding down column~$2$:
  \[
    \det(A-\lambda I)=(2-\lambda)\det\begin{bmatrix}3-\lambda&1\\0&1-\lambda\end{bmatrix}=(2-\lambda)(3-\lambda)(1-\lambda).
  \]
  The eigenvalues are $1,2,3$.
  \ii For $\lambda=1$:
  \[
    A-I=\begin{bmatrix}2&0&1\\-1&1&-1\\0&0&0\end{bmatrix}.
  \]
  Row~$1$ gives $x_3=-2x_1$ and Row~$2$ gives $x_2=x_1+x_3=-x_1$. With $x_1$ free, a basis of the eigenspace is $\{(1,-1,-2)\}$. (Check: $A(1,-1,-2)^{\mathsf T}=(1,-1,-2)^{\mathsf T}$.)
\end{steps}
\end{prob}

\begin{prob}\textbf{\xlink{math54}{2026-10-05-coordinates-eigenvalues}{prob-P6}{P6}.}\par\idea{Use $\lambda^2-(\tr A)\lambda+\det A$ for $2\times 2$ matrices.}
\begin{steps}
  \ii First matrix: $\tr=4$, $\det=4-49=-45$, so $\lambda^2-4\lambda-45=(\lambda-9)(\lambda+5)$, with eigenvalues $9,-5$ (check: $9+(-5)=4$, $9(-5)=-45$).
  \ii Second matrix: $\tr=8$, $\det=15-12=3$, so characteristic polynomial $\lambda^2-8\lambda+3$, with eigenvalues $\lambda=\frac{8\pm\sqrt{64-12}}{2}=4\pm\sqrt{13}$.
\end{steps}
\end{prob}

\begin{prob}\textbf{\xlink{math54}{2026-10-05-coordinates-eigenvalues}{prob-P7}{P7}.}\par\idea{The matrix is upper triangular, so eigenvalues sit on the diagonal (\xlink{math54}{2026-10-05-coordinates-eigenvalues}{thm-lay-5-1}{Thm.~5.1}).}
\begin{steps}
  \ii Characteristic polynomial: $(2-\lambda)^2(-1-\lambda)(0-\lambda)$.
  \ii Eigenvalues: $2$ (multiplicity $2$), $-1$ (multiplicity $1$), $0$ (multiplicity $1$).
  \ii Because $0$ is an eigenvalue, the matrix is \textbf{not invertible}.
\end{steps}
\end{prob}

\begin{prob}\textbf{\xlink{math54}{2026-10-05-coordinates-eigenvalues}{prob-P8}{P8}.}\par\idea{Apply $A$ twice to $A\vv x=\lambda\vv x$; for (b) multiply on the left by $A^{-1}$.}
\begin{steps}
  \ii \textbf{(a)} $A^2\vv x=A(A\vv x)=A(\lambda\vv x)=\lambda A\vv x=\lambda^2\vv x$, and $(A+3I)\vv x=A\vv x+3\vv x=(\lambda+3)\vv x$. Since $\vv x\ne\vv 0$, $\vv x$ is an eigenvector of $A^2$ with eigenvalue $\lambda^2$ and of $A+3I$ with eigenvalue $\lambda+3$.
  \ii \textbf{(b)} If $\lambda=0$, then $A\vv x=\vv 0$ with $\vv x\ne\vv 0$, contradicting invertibility of $A$, so $\lambda\ne 0$. Multiplying $A\vv x=\lambda\vv x$ on the left by $A^{-1}$ gives $\vv x=\lambda A^{-1}\vv x$, and dividing by $\lambda$ yields $A^{-1}\vv x=\lambda^{-1}\vv x$.
\end{steps}
\end{prob}

\end{document}
